\section{Giới hạn và việc còn lại}

\subsection{Những gì chưa được phép kết luận}

\paragraph{Phase 1.}
Không phân định được mô hình dense nào tốt nhất --- các khoảng tin cậy chồng
lấn. Nghiêm trọng hơn, nhánh dense \textbf{chưa tái lập được} giữa các lần
chạy: độ trôi đo được $0{,}0143$ nDCG@5, \emph{lớn hơn} khoảng cách $0{,}0094$
giữa hai mô hình dense đang so sánh. Nguyên nhân đã xác định được trong mã
nguồn: không đặt seed cho HNSW, không cố định \texttt{batch\_size} khi tính
embedding, không gọi \texttt{torch.manual\_seed}. Chừng nào chưa sửa, mọi xếp
hạng nội bộ giữa các mô hình dense đều không có ý nghĩa. Kết quả chunk size là
null; hybrid là ``không chứng minh được có lợi'', không phải ``đã chứng minh có
hại''.

\paragraph{Phase 2.}
Answer Correctness \NumHoAC{} là một \textbf{cận dưới}. EDA ước lượng
\NumNoiseFloor{}--\NumNoiseCeiling{} số câu mang nhãn có vấn đề, và audit tay
thấy 23/30 câu điểm thấp thực ra đúng về ngữ nghĩa. Con số thực nằm đâu đó
trong dải giữa \NumHoAC{} và một trần cao hơn đáng kể mà thí nghiệm này không
đo được.

Judge cũng là một LLM. Nó đã được tách sang nhà cung cấp khác generator --- một
sửa lỗi trực tiếp từ phát hiện của EDA rằng notebook cũ chấm bài bằng chính mô
hình sinh --- nhưng \textbf{chưa được hiệu chuẩn với nhãn người}. Tập 20 câu
dành cho việc hiệu chuẩn đã chuẩn bị nhưng chưa dùng.

Answer Relevancy giảm từ \NumDevAnsRel{} xuống \NumHoAnsRel{}. Đây là hệ quả có
chủ đích của P2, không phải suy thoái: metric ấy thưởng câu trả lời đầy đủ ngữ
cảnh, còn dữ liệu này thưởng câu trả lời cụt. Nêu ra vì nó là một đánh đổi thật,
không phải để giấu.

Cuối cùng, giải đấu phân tầng ở cả hai phase \textbf{không phát hiện được tương
tác} giữa các lựa chọn ở những vòng khác nhau.

\subsection{Việc còn lại}

\begin{center}\small
\begin{tabularx}{0.94\textwidth}{@{}p{0.7cm} L p{1.2cm} p{3.2cm}@{}}
\toprule
\textbf{\#} & \textbf{Việc} & \textbf{Phase} & \textbf{Ghi chú} \\
\midrule
1 & Audit mù 30 cặp P0 $\leftrightarrow$ winner, hai người chấm độc lập & 2 & Cần chia việc giữa hai người \\
2 & Lặp 25 câu $\times$ 2 cấu hình để đo độ ổn định API & 2 & \\
3 & Hiệu chuẩn judge trên 20 câu đã dành sẵn & 2 & Tập đã chuẩn bị \\
4 & Chạy held-out trên 150 bài final-test & 1 & Nghiệm thu Phase 1 \\
5 & Chạy lại vòng 1 trên chỉ mục đã đóng gói & 1 & Cố định lại số dense \\
6 & Bốn biểu đồ 300 DPI theo test plan Phase 1 §6 & 1 & Chưa có cái nào \\
7 & Duyệt tay dataset abstention 200 case & 3 & 3 loại cần reviewer thứ hai \\
\bottomrule
\end{tabularx}
\end{center}

Phase 3 --- dạy hệ thống biết từ chối trả lời khi bằng chứng không đủ --- là
bước tiếp theo, và phân tầng ở mục 4 cho thấy vì sao nó đáng làm:
\NumHoNMiss{}/\NumNHeldout{} câu held-out không có bằng chứng nào để trả lời
đúng, nhưng hệ thống vẫn trả lời chúng.

\subsection{Tổng kết}

\begin{enumerate}\small
  \item \textbf{Sparse thắng dense} (\NumSparseVsDense{} nDCG@5,
        \NumSparseVsDenseCI{}) vì câu hỏi NewsQA neo vào tên riêng, ngày tháng
        và con số. EDA dự báo trước, thực nghiệm xác nhận.
  \item \textbf{Prompt đúng \emph{dạng} là đòn bẩy lớn nhất ở tầng sinh} ---
        không phải vì kỹ thuật viết prompt khéo, mà vì nó khớp với hình dạng
        của nhãn.
  \item \textbf{Guardrail đăng ký trước đã loại một cấu hình có điểm cao hơn},
        và đó chính xác là việc nó sinh ra để làm.
  \item \textbf{Prompt không bị overfit}: khoảng cách development $\rightarrow$
        held-out quy về việc truy xuất khó hơn, không phải prompt kém tổng quát.
\end{enumerate}

Số công bố: Answer Correctness \textbf{\NumHoAC{}} trên \NumNHeldout{} câu
held-out, cấu hình \NumWinner{}, chạy đúng một lần.

\subsection*{Tài liệu và mã nguồn liên quan}

\begin{itemize}\small
  \item Prompt nguyên văn và cách một request được ghép: \texttt{docs/prompts.md}
  \item Báo cáo chi tiết từng phase: \texttt{docs/reports/phase1/report.md},
        \texttt{docs/reports/phase2/report.md}
  \item Kế hoạch đăng ký trước: \texttt{docs/Detailed Test Plans/}
  \item Quyết định winner và bằng chứng held-out:
        \texttt{docs/reports/phase2/phase2b\_winner\_decision.json},
        \texttt{docs/reports/phase2/heldout/}
  \item Sinh lại toàn bộ con số trong tài liệu này:
        \texttt{python scripts/build\_latex\_numbers.py}
\end{itemize}
